\section{Method}
\label{sec:method}

% 2026-10-14: retain the complete-contribution and receiving-operator definitions from Section 39.
% Collaboration scoring, strength selection, and dependent allocation rules are pending redesign.
% Sources: ../FedLLM/ACL_Logs/method-39_progressive_recipient_allocation.md;
% ../FedLLM/exp/_archive_exps/1003_acl_progressive_exchange/design/progressive_allocation_report.md;
% ../FedLLM/alg/acl_1003_progressive_exchange.py and ../FedLLM/utils/progressive_{allocation_utils,metric}.py.
% Paper s_i denotes receiving capacity; the implementation calls it q_i.
% Writing references: ref/FedMosaic_SIGIR26_CR and ref/DarkDistill_KDD25_CR; see WRITING_STYLE.md.

\TODO{Section 4.1 Overview}

\TODO{Section 4.2 Allocation}
\TODO{Take collboaration weight into method design}

\TODO{Section 4.3 Adaptive Collaboration}
\TODO{Take allocation result into method design}

\TODO{Section 4.4 Putting it Together}

\subsection{Overview}
\label{sec:method:overview}

\fakeparagraph{Key Idea}
Under rank heterogeneity, a client should not have to choose between exposing its
entire adapter and receiving no external knowledge.
\sysname therefore decouples contribution from reception.
Each client contributes its complete locally trained adapter to a round-level pool,
while only a client-specific portion of its fixed rank budget is exposed to
incoming updates.
For client $i$, the receiving capacity $s_i$ and protected capacity $p_i=r_i-s_i$
determine where external knowledge may enter, whereas a client-specific
collaboration strength determines how much of that knowledge is applied.
This separation preserves the client's total rank budget and local adaptation
capacity while allowing useful cross-client transfer.

\fakeparagraph{Challenge-to-Component Map}
This design follows the two challenges in \secref{sec:problem:challenges}.
The allocation module determines how client $i$ divides its fixed budget $r_i$
between receiving and protected capacity, rather than imposing one sharing width
on all clients.
The collaboration module then determines whether the available external influence
benefits client $i$ and, if so, how strongly it should be applied.
The final putting-it-together procedure coordinates these decisions with complete
contribution, so restricting a client's reception does not remove its knowledge
from the pool available to other clients.

\fakeparagraph{Workflow}
At the beginning of each round, clients locally train both portions of their
adapters under the fixed budgets in \equref{equ:heterogeneous-objective} and
contribute their complete effective updates.
The server constructs the external target for client $i$ from the other clients'
updates using fixed training-size weights, as in \equref{equ:acl-target}.
The allocation module selects a receiving capacity $s_i$ and thereby protects
$p_i=r_i-s_i$ components.
Given this capacity, the collaboration module constructs the residual between
the external target and client $i$'s local update in the permitted subspace and
selects its reception strength, including zero when external influence is not
useful.
The final update is applied only to the receiving portion, producing the
personalized adapter used in the next round; the protected portion remains
untouched by direct incoming modification.

\fakeparagraph{Architecture}
\begin{itemize}
    \item \textbf{Allocation Module (\secref{sec:method:allocation}).}
    It determines the client-specific capacity $s_i$ within the fixed rank
    budget and reserves the remaining capacity for local adaptation.
    \item \textbf{Collaboration Module (\secref{sec:method:collaboration}).}
    It computes the external influence available within the selected receiving
    subspace and determines its client-specific reception strength.
    \item \textbf{Putting It Together.}
    It orders complete contribution, capacity allocation, and selective reception
    into one round-level procedure, while freezing the contribution pool before
    applying incoming updates.
\end{itemize}

Together, these components turn selective rank sharing into a coordinated
allocation-and-collaboration process: allocation controls where external
knowledge may enter, collaboration controls whether it should enter, and the
round procedure preserves complete contribution across clients.
The following subsections detail these components in the same order as the
workflow.

\subsection{Local Training and Complete Contribution}
\label{sec:method:contribution}

For the proposed client-specific allocation, client $i$ partitions its adapter as
\begin{equation}
    B_iA_i=B_{i,s}A_{i,s}+B_{i,p}A_{i,p},
    \qquad p_i=r_i-s_i,\quad s_i\in\{0,\ldots,r_i\}.
    \label{equ:capacity-partition}
\end{equation}
Client $i$ minimizes its local term in \equref{equ:heterogeneous-objective} from its preceding personalized state, training both portions in \equref{equ:capacity-partition} with total width $r_i$ and LoRA scaling fixed.
It then uploads the complete adapter, including the protected components, so that reducing its reception permission does not remove its knowledge from the contribution pool.
Within a round, $B_i,A_i$ denote the factors after local training, with scaling absorbed as in \secref{sec:problem}.
The server pools the other clients' complete effective updates into
\begin{equation}
    M_{-i}=\sum_{j\ne i}w_{ij}B_jA_j,
    \qquad
    w_{ij}=\frac{|\mathcal{D}_j|}{\sum_{k\ne i}|\mathcal{D}_k|},
    \label{equ:acl-target}
\end{equation}
where $|\mathcal{D}_j|$ counts local training examples, excluding selection data.
This forms an all-to-all contribution pool with every other client as an available source and training-size weights for the pooled target.
The donor proportions are fixed across candidate widths, leaving the receiving subspace and strength to determine how the pooled knowledge affects each client.
The server freezes all uploads before constructing receiving directions, so a change made during reception enters other clients' targets only in the next round.
Without an external contribution, the client retains its local state.

\subsection{Reception at a Given Capacity}
\label{sec:method:collaboration}

\fakeparagraph{Constructing the Receiving Direction}
For a candidate width $s$, we permit reception into the first $s$ native components and protect the remaining $r_i-s$ components.
Component identities retain their native order across rounds.
Since equal component indices need not represent the same directions across clients, we compare their complete effective updates and project the difference into client $i$'s permitted subspace.
For $s>0$, the server computes
\begin{equation}
    \Delta B_{i,s}=(M_{-i}-B_iA_i)A_{i,s}^{\dagger},
    \label{equ:acl-direction}
\end{equation}
where $A_{i,s}^{\dagger}$ is the Moore--Penrose pseudoinverse.
With strength $\lambda\in[0,1]$, reception gives
\begin{equation}
    \begin{aligned}
        B_i(s,\lambda)&=[\,B_{i,s}+\lambda\Delta B_{i,s},\ B_{i,p}\,],
        \qquad A_i(s,\lambda)=A_i,\\
        B_i(s,\lambda)A_i
        &=B_iA_i+\lambda(M_{-i}-B_iA_i)A_{i,s}^{\dagger}A_{i,s}.
    \end{aligned}
    \label{equ:acl-receive}
\end{equation}
The projection transfers the part of the residual represented in the selected row space, while changing only the receiving $B$ columns.
Subtracting the complete local update accounts for content already retained in the protected block; the direction vanishes when the complete local and external updates coincide.
The adapter retains $r_i$ component slots and algebraic rank at most $r_i$.
Zero width or zero strength preserves the original factors exactly.

\fakeparagraph{Collaboration Metric}
\TODO{Define a computable collaboration metric and its use in selecting reception strength, including zero reception, after the design and corresponding experiments are ready.}

\subsection{Progressive Capacity Allocation}
\label{sec:method:allocation}

\TODO{Reformulate capacity comparisons and update rules around the new collaboration metric once it is defined and evaluated.}
